\chapter{2009年广东卷（理）}

\section{单选题}

\begin{problem}
已知全集 \(U = \R\)，集合 \(M = \{x \mid -2 \leqslant x - 1 \leqslant 2\}\) 和 \(N = \{x \mid x = 2k - 1, k = 1, 2, \cdots\}\) 的关系的韦恩 (Venn) 图如图所示，则阴影部分所集合的元素共有（\quad）

\choices
    {3 个}
    {2 个}
    {1 个}
    {无穷个}
\end{problem}

\begin{problem}
设 \(z\) 是复数，\(\alpha(z)\) 表示满足 \(z^n = 1\) 的最小正整数 \(n\)，则对虚数单位 \(\i\)，\(\alpha(\i) =\)（\quad）

\choices
    {8}
    {6}
    {4}
    {2}
\end{problem}

\begin{problem}
若函数 \( y = f(x) \) 是函数 \( y = a^x (a > 0, \text{且 } a \neq 1) \) 的反函数，其图象经过点 \( (\sqrt{a}, a) \)，则 \( f(x) = \)（\quad）

\choices
    {\(\log_2 x\)}
    {\(\log_{\frac{1}{2}}x\)}
    {\(\frac{1}{2^\e}\)}
    {\(x^{2}\)}
\end{problem}

\begin{problem}
已知等比数列 \(\{a_{n}\}\) 满足 \(a_{n} > 0, n = 1,2,\dots\)，且 \(a_{2} \cdot a_{2n-3} = 2^{2n} (n \geqslant 3)\)，则当 \(n \geqslant 1\) 时，\(\log_2 a_1 + \log_2 a_3 + \dots + \log_2 a_{2n-1} =\)（\quad）

\choices
    {\(n(2n - 1)\)}
    {\((n + 1)^{2}\)}
    {\(n^{2}\)}
    {\((n - 1)^{2}\)}
\end{problem}

\begin{problem}
给定下列四个命题: \circled{1} 若一个平面内的两条直线与另一个平面都平行，那么这两个平面相互平行; \circled{2} 若一个平面经过另一个平面的垂线，那么这两个平面相互垂直; \circled{3}垂直于同一直线的两条直线相互平行； \circled{4} 若两个平面垂直，那么一个平面与它们的交线不垂直的直线与另一个平面也不垂直. 其中，为真命题的是（\quad）

\choices
    {\circled{1}和\circled{2}}
    {\circled{2}和\circled{3}}
    {\circled{3}和\circled{4}}
    {\circled{2}和\circled{4}}
\end{problem}

\begin{problem}
一质点受到平面上的三个力 \(F_{1}, F_{2}, F_{3}\) (单位: 牛顿) 的作用而处于平衡状态. 已知 \(F_{1}: F_{3}\) 成 \(60^{\circ}\) 角，且 \(F_{1}, F_{2}\) 的大小分别为 2 和 4，则 \(F_{3}\) 的大小为（\quad）

\choices
    {6}
    {2}
    {\(2\sqrt{5}\)}
    {\(2\sqrt{7}\)}
\end{problem}

\begin{problem}
2010 年广州亚运会组委会要从小张、小赵、小李、小罗、小王五名志愿者中选派四人分别从事翻译、导游、礼仪、司机四项不同工作，若其中小张和小赵只能从事前两项工作，其余三人均能从事这四项工作，则不同的选派方案共有（\quad）

\choices
    {36 种}
    {12 种}
    {18 种}
    {48 种}
\end{problem}

\begin{problem}
已知甲、乙两车由同一起点同时出发，并沿同一路线 (假定为直线) 行驶. 甲车、乙车的速度曲线分别为 \(v_{\text{甲}}\) 和 \(v_{\text{乙}}\) (如图所示). 那么对于图中给定的 \(t_0\) 和 \(t_1\)，下列判断中一定正确的是（\quad）

\choices
    {在 \( t_{1} \) 时刻，甲车在乙车前面}
    {\(t_1\) 时刻后，甲车在乙车后面}
    {在 \( t_{0} \) 时刻，两车的位置相同}
    {\(t_0\) 时刻后，乙车在甲车前面}
\end{problem}

\section{填空题}

\begin{problem}
随机抽取某产品 \(n\) 件，测得其长度分别为 \(a_1, a_2, \cdots, a_n\)，则如图所示的程序框图输出的 \(s = \fillinblank{}\fillinblank{}, s\) 表示的样本的数字特征是 \(\fillinblank{}\fillinblank{}\). (注: 框图中的赋值符号“=”也可以写成“←”或“:=”)
\end{problem}

\begin{problem}
若平面向量 \(a, b\) 满足 \(|a + b| = 1, a + b\) 平行于 \(x\) 轴，\(b = (2, -1)\)，则 \(a = \)  \fillinblank{}.
\end{problem}

\begin{problem}
已知椭圆 \(G\) 的中心在坐标原点，长轴在 \(x\) 轴上，离心率为 \(\frac{\sqrt{3}}{2}\)，且 \(G\) 上一点到 \(G\) 的两个焦点的距离之和为 12 ，则椭圆 \(G\) 的方程为 \fillinblank{} \fillinblank{}.
\end{problem}

\begin{problem}
已知离散型随机变量 X 的分布列如下表. 若 EX = 0, DX = 1，则 a = ，b =\fillinblank{}. X -1 0 1 2 P a b c \( \frac{1}{12} \)
\end{problem}

\begin{problem}
若直线 \( l_{1}:\left\{\begin{aligned}x&=1-2t,\\ y&=2+kt,\end{aligned}\right. \) (t 为参数) 与直线 \( l_{2}:\left\{\begin{aligned}x&=s,\\ y&=1-2s,\end{aligned}\right. \) (s 为参数) 垂直，则 \( k=\frac{|s|}{|x+2|}\geqslant1 \) 的实数解为 \fillinblank{} \fillinblank{}.
\end{problem}

\begin{problem}
如图，点 \(A, B, C\) 是圆 \(O\) 上的点，且 \(AB = 4\)，\(\angle ACB = 45^{\circ}\)，则圆 \(O\) 的面积等于 \fillinblank{} \fillinblank{}.
\end{problem}

\section{解答题}

\begin{problem}
已知向量 \( a = (\sin \theta, -2) \) 与 \( b = (1, \cos \theta) \) 互相垂直，其中 \( \theta \in \left(0, \frac{\pi}{2}\right) \) . (1) 求 \( \sin\theta\cos\theta \) 的值. (2) 若 \( \sin(\theta-\varphi)=\frac{\sqrt{10}}{10}, 0<\varphi<\frac{\pi}{2} \) ，求 \( \cos\varphi \) 的值.
\end{problem}

\begin{problem}
根据空气质量指数 API（为整数）的不同，可将空气质量分级如下表：对某城市一年（365 天）的空气质量进行监测，获得 API 数据按照区间 [0, 50]，(50, 100), (100, 150], (150, 200], (200, 250], (250, 300] 进行分组，得到频率分布直方图如图. API 0 \&lt; t \&lt; 50 1\textasciitilde{}50 101\textasciitilde{}150 151\textasciitilde{}250 250\textasciitilde{}350 350\textasciitilde{}475 475\textasciitilde{}650 \&gt;650 \&gt;650 区间 无 无 0.5 0.5 0.5 0.5 0.5 0.5 0.5 (1) 求直方图中 x 的值; (2) 计算一年中空气质量分别为良和轻微污染的天数. (3) 求流域城市一期至少有 2 天的空气质量均为优良或轻微污染的概率. (结果用分数表示，已知 \( 5^{x} = 7812^{y}, 2^{x} = 128, \frac{1}{1825}, \frac{1}{365}, \frac{1}{1825}, \frac{3}{1825} + \frac{3}{1825} + \frac{8}{1825} = \frac{123}{1435}, 365 \times 73 \times 5) \)
\end{problem}

\begin{problem}
如图，已知正方体 \(ABCD - A_{1}B_{1}C_{1}D_{1}\) 的棱长为2，点 \(E\) 是正方形 \(BCC_{1}B_{1}\) 的中心，点 \(F\) 、 \(G\) 分别是棱体 \(C_1D_1\) ， \(A_{1}A_{1}\) 的中点.设点 \(E_{1}\) ， \(C_{1}\) 分别是点 \(E\) 、 \(G\) 在平面 \(DCC_{1}D_{1}\) 内的正投影. (1) 求以 E 为顶点，以四边形 FGAE 在平面 \( DCC_{1}D_{1} \) 内的正投影为底面边界的棱锥的体积; (2) 证明: 直线 \( FG_{1} \perp FEE_{1} \) ; (3) 求异面直线 \( E_{1}G_{1} \) 与 EA 所成角的正弦值.
\end{problem}

\begin{problem}
已知二次函数 \( y = g(x) \) 的导函数的图象与直线 \( y = 2x \) 平行，且 \( y = g(x) \) 在 \( x = -1 \) 处取得极小值 \( m - 1 (m \neq 0) \). 设 \( f(x) = \frac{g(x)}{x} \). (1) 若曲线 \( y = f(x) \) 上的点 \( P \) 到点 \( Q(0,2) \) 的距离的最小值为 \( \sqrt{2} \)，求 \( m \) 的值； (2) \(k(k \in \R)\) 如为取值时，函数 \(y = f(x) - kx\) 存在零点，并求出零点.
\end{problem}

\begin{problem}
已知曲线 \(C_x: x^2 - 2nx + y^2 = 0 (n = 1, 2, \cdots)\). 从点 \(P(-1,0)\) 向曲线 \(C_x\) 引斜率为 \(k_0 (k_n > 0)\) 的切线 \(L_0\)，切点为 \(P_n(x_n, y_n)\). (1) 求数列 \( \{x_{n}\} \) 与 \( \{y_{n}\} \) 的通项公式； (2) 证明: \(x_{1} \cdot x_{3} \cdot x_{5} \cdots \cdot x_{2n-1} < \sqrt{\frac{1 - x_{n}}{1 + x_{n}}} < \sqrt{2} \sin \frac{x_{n}}{y_{n}}\).
\end{problem}

\begin{problem}
已知曲线 \(C: y = x^2\) 与直线 \(l: x - y + 2 = 0\) 交于两点 \(A(x_A, y_A)\) 和 \(B(x_B, y_B)\)，且 \(x_A < x_B\). 记曲线 \(C\) 在点 \(A\) 和点 \(B\) 之间那一段 \(L\) 与线段 \(AB\) 所围成的平面区域 (含边界) 为 \(D\). 设点 \(P(s, t)\) 是 \(L\) 上的任一点，且点 \(P\) 与点 \(A\) 和点 \(B\) 均不重合. (1) 若点 \(Q\) 是线段 \(AB\) 的中点，试求线段 \(PQ\) 的中点 \(M\) 的轨迹方程; (2) 若曲线 \(G: x^{2} - 2ax + y^{2} - 4y + a^{2} + \frac{51}{25} = 0\) 与 \(D\) 有公共点，试求 \(a\) 的最小值.
\end{problem}

